% ================================================================
%  第 3 章  弱收敛、紧性与应用
% ================================================================
\chapter{弱收敛、紧性与应用}
\label{chap:3}

本章中, 我们将对子流形引入一种弱收敛的概念, 称为变分流形收敛: 这里把
子流形等同于一个可以对其积分的测度. 我们将看到, 极小子流形有一个自然的
推广, 并且经典理论的大部分都可以延拓到这一框架. 接下来我们将证明
Michael--Simon Sobolev 不等式. 然后我们将证明 \cite{CM16} 中给出的、
对前几章所讨论的极小曲面的经典 Bernstein 定理的一个推广.

各种 Bernstein 定理都表明: 在维数不超过七时, 面积最小化超曲面都必须是
仿射的. Allard 正则性定理 \cite{Al1} 给出了一个在一切维数下都成立的较弱
形式. 即: 存在
\begin{equation}
\delta=\delta(k,n)>0
\label{eq:3.1}
\end{equation}
使得, 若 $\Sigma^k\subset\R^n$ 是 $k$ 维完全浸入极小子流形, 且对一切 $x$
与 $r$ 都有
\begin{equation}
\Vol(B_r(x)\cap\Sigma)\leq(1+\delta)\Vol(B_r\subset\R^k),
\label{eq:3.2}
\end{equation}
则 $\Sigma$ 是一张仿射 $k$ 维平面. 更一般地, 即使 $\Sigma$ 带有奇点, 该
定理仍然成立, 例如当 $\Sigma$ 是驻定整 $k$-变分流形时(见下面的\defref{def:3.3}、
3.4 与 3.10).

本章中我们将证明: 事实上, 密度的上界(变分流形密度的定义见 \eqref{eq:3.34})给出
了包含该极小曲面的最小仿射子空间维数的一个上界. 我们将从欧氏空间中任意
余维的 $k$-可求长驻定变分流形上坐标函数的性质(事实上更一般地, 从调和函数
的性质)出发来推出这个定理. 进一步的讨论与更一般的结果可参看原始论文
\cite{CM16}.

最后, 在本章的最后一节中, 我们将讨论另一种弱收敛的概念, 称为气泡收敛,
它与 Sacks 与 Uhlenbeck 以及 Gromov 的工作有关. 随后我们将证明: 气泡
收敛蕴含变分流形收敛.

% ================================================================
\section{变分流形理论}
\label{sec:3.1}

我们首先简要叙述变分流形的基本理论. 为记号方便, 我们把讨论限制在 $\R^n$
中的变分流形; 无须太多额外工作就可以考虑光滑 $n$ 维流形中的变分流形.
应当把变分流形看作子流形的一种推广. 这一类广义子流形足够宽, 足以给出
一般的紧性结果并容许奇点, 但又足够窄, 仍然具有几何意义. 例如, 我们可以
为变分流形定义第一变分.

在变分流形的理论中, 曲面被等同于某些 \newterm{radon-measure}{Radon 测度}{Radon measure}. 基本的紧性结果于是
可以由 Radon 测度的紧性定理推出.

\begin{definition}
\label{def:3.1}
设 $X$ 是局部紧的可分空间. 若 Borel 正则测度 $\mu$ 在 $X$ 的紧子集上
有限, 则称 $\mu$ 为 \textbf{Radon 测度}.
\end{definition}

下述标准紧性结果是把子流形看作 Radon 测度的主要好处:

\begin{theorem}[Radon 测度的紧性]
\label{thm:3.2}
设 $X$ 是局部紧的可分空间, $\mu_j$ 是 $X$ 上的一列 Radon 测度, 使得对
任意闭包紧的开集 $U$ 都有
\begin{equation}
\sup_j\mu_j(U)<\infty.
\label{eq:3.3}
\end{equation}
则存在 Radon 测度 $\mu$ 与子列 $\mu_{j'}$, 使得 $\mu_{j'}\to\mu$; 即:
若 $f$ 是 $X$ 上具有紧支集的连续函数, 则
\begin{equation}
\lim_{j'\to\infty}\int_X f\,d\mu_{j'}=\int_X f\,d\mu.
\label{eq:3.4}
\end{equation}
\end{theorem}

设 $\Sigma^k\subset\R^n$ 是光滑的真嵌入子流形. 有一种明显的方式给
$\Sigma$ 配一个 Radon 测度 $\mu$; 即令
\begin{equation}
\mu(A)=\Vol(A\cap\Sigma)
\label{eq:3.5}
\end{equation}
对 Borel 集 $A\subset\R^n$ 成立. 只要这些子流形具有一致的局部面积界,
\thmref{thm:3.2} 就给出子流形(或相应的测度)的紧性. 即: 设 $C<\infty$,
$\Sigma_j\subset\R^n$ 是一列 $k$ 维子流形, 使得对一切 $x\in\R^n$ 与一切
$j$ 都有
\begin{equation}
\Vol(B_1(x)\cap\Sigma_j)<C.
\label{eq:3.6}
\end{equation}
按 \eqref{eq:3.5} 定义相应的 Radon 测度 $\mu_j$. 由\thmref{thm:3.2}, 存在子列 $j'$ 与
Radon 测度 $\mu$, 使得 $\mu_{j'}$ 弱收敛于 $\mu$.

把一列子流形看作测度所要付出的代价是: 极限仅仅是一个测度(未必表现出
任何几何). 特别地, 如果我们试图从一个极小化序列中提取面积最小化的
子流形, 那么就无法断言极限是极小的(即第一变分为零).

设 $G(k,n)$ 是 $\R^n$ 中过原点的(无向)$k$ 维平面构成的空间(于是
$G(n-1,n)$ 是射影 $(n-1)$ 空间). 令
\begin{equation}
\pi:\R^n\times G(k,n)\to\R^n
\label{eq:3.7}
\end{equation}
表示投影 $\pi(x,\omega)=x$. 注意 $\pi$ 是固有映射. 我们还要对 $x\in\R^n$
处的 $k$ 维平面 $\omega$ 定义散度与梯度如下: 若 $X$ 是 $C^1$ 向量场, 则
定义它关于 $(x,\omega)$ 的散度为
\begin{equation}
\dive_\omega X=\sum_{i=1}^k\langle E_i,\nabla_{E_i}X\rangle,
\label{eq:3.8}
\end{equation}
其中 $E_i$ 是 $\omega$ 的标准正交基, $\nabla$ 是欧氏导数. 同样, 若 $u$ 是
$C^1$ 函数, 则令
\begin{equation}
\nabla_\omega u=\sum_{i=1}^k E_i(u)\,E_i.
\label{eq:3.9}
\end{equation}

\begin{definition}[变分流形]
\label{def:3.3}
$\R^n$ 上的 \textbf{$k$-变分流形} $T$ 是 $\R^n\times G(k,n)$ 上的
Radon 测度. $T$ 的\textbf{权} $\mu_T$ 是 $\R^n$ 上由
$\mu_T(E)=T(\pi^{-1}E)$ 给出的 Radon 测度, $T$ 的支集就是 $\mu_T$ 的
支集, 而 $T$ 在集合 $U\subset\R^n$ 上的\textbf{质量}就是 $\mu_T(U)$.
\end{definition}

再设 $\Sigma^k\subset\R^n$ 是光滑的真嵌入子流形. 定义
$T\Sigma\subset\R^n\times G(k,n)$ 为: $(x,\omega)\in T\Sigma$ 当且仅当
$\omega$ 是 $\Sigma$ 在 $x$ 处的切 $k$ 维平面. $\Sigma$ 给出一个
$k$-变分流形 $T$, 它对 Borel 集
$A\subset\R^n\times G(k,n)$ 定义为
\begin{equation}
T(A)=\Vol(\pi(A\cap T\Sigma)),
\label{eq:3.10}
\end{equation}
这里 $\Vol$ 是 $k$ 维体积. 此时, $T$ 在集合 $U$ 上的质量就是 $\Sigma$ 在
$U$ 中的体积.

有了这个定义, 就很容易看出变分流形在映射下的行为. 即: 若
$F:\R^n\to\R^n$ 是 $C^2$ 映射(于是 $F$ 与 $dF$ 都是 $C^1$), 则 $F$
诱导一个 $C^1$ 映射 $\mathcal{F}$, 从 $\R^n\times G(k,n)$ 到
$\R^n\times\bigcup_{\ell\leq k}G(\ell,n)$:
\begin{equation}
\mathcal{F}(x,\omega)=(F(x),dF_x(\omega)).
\label{eq:3.11}
\end{equation}
若 $T$ 是 $k$-变分流形, 且 $F$ 在 $T$ 的支集上是固有的, 则定义(带权的)
\textbf{前推}变分流形 $F(T)$ 为
\begin{equation}
F(T)(A)=\int_{\mathcal{F}^{-1}A}J_\omega F(x)\,dT(x,\omega),
\label{eq:3.12}
\end{equation}
其中 $J_\omega F(x)$ 是由
\begin{equation}
J_\omega F(x)=\bigl(\det(dF_x|_\omega)^t\circ(dF_x|_\omega)\bigr)^{\frac12}
\label{eq:3.13}
\end{equation}
给出的 Jacobi 因子. 被积式 \eqref{eq:3.13} 在 $\mathcal{F}^{-1}$ 无定义的集合上
为零, 因此由 \eqref{eq:3.12} 给出的前推变分流形 $F(T)(A)$ 是定义良好的. 当 $T$
来自光滑曲面 $\Sigma$ 时, $J_\omega F(x)$ 正是拉回度量的行列式的平方根,
此时 \eqref{eq:3.12} 就是通常的面积公式(例如见 \cite{Fe}).

利用变分流形的这一变换规则, 我们可以计算变分流形``质量''的第一变分.
设
\begin{equation}
F(x,t):\R^n\times\R\to\R^n
\label{eq:3.14}
\end{equation}
是一族 $C^2$ 微分同胚, 使得 $F(\cdot,0)$ 是恒等映射, 且 $F(\cdot,t)$ 在
一个固定的紧集 $K$ 之外是恒等映射. 记 $F_t(x)$ 为向量场
$\frac{\partial F}{\partial t}(x,0)$(于是 $F_t$ 在紧集 $K$ 之外为零), 并
设 $K\subset U\subset\R^n$. 在 \eqref{eq:3.12} 的积分号下求导, 并完全照第 1 章的
论证, 得到
\begin{equation}
\Bigl.\frac{d}{dt}\Bigr|_{t=0}\mu_{F(T,t)}(U)
 =\int_{U\times G(k,n)}\dive_\omega F_t(x)\,dT(x,\omega).
\label{eq:3.15}
\end{equation}

若变分流形对任意紧支集 $C^1$ 向量场的第一变分都为零, 则称它是
\textbf{驻定}的.

\begin{definition}
\label{def:3.4}
若对任意紧支集 $C^1$ 向量场 $X$ 都有
\begin{equation}
\int_{\R^n\times G(k,n)}\dive_\omega X\,dT(x,\omega)=0,
\label{eq:3.16}
\end{equation}
则称变分流形 $T$ 是\newterm{stationary}{驻定}{stationary}的.
\end{definition}

注意, 由 \eqref{eq:3.15}, 这个定义等价于要求变分流形在一切紧支集 $C^2$ 变分下
都是质量泛函的临界点. 还要注意, 若变分流形来自光滑子流形, 则第一变分
公式表明: 它是驻定的, 当且仅当它的平均曲率为零.

这个定义的一个重要推论是下述对前面讨论的变分流形紧性定理的补充:

\begin{proposition}
\label{prop:3.5}
若驻定变分流形序列 $T_j$ 弱收敛于变分流形 $T$, 则 $T$ 也是驻定的.
\end{proposition}

\begin{proof}
给定具有紧支集的 $C^1$ 向量场 $X$, 我们可以用
\begin{equation}
f(x,\omega)=\dive_\omega X
\label{eq:3.17}
\end{equation}
定义一个具有紧支集的连续函数 $f$. 由于 $T_j$ 驻定,
\begin{equation}
\int_{\R^n\times G(k,n)}f(x,\omega)\,dT_j(x,\omega)
 =\int_{\R^n\times G(k,n)}\dive_\omega X\,dT_j(x,\omega)=0.
\label{eq:3.18}
\end{equation}
于是由变分流形收敛性以及 $f$ 在 $\R^n\times G(k,n)$ 上的紧支性, 得到
\begin{equation}
\begin{aligned}
\int_{\R^n\times G(k,n)}\dive_\omega X\,dT(x,\omega)
&=\int_{\R^n\times G(k,n)}f\,dT(x,\omega)\\
&=\lim_{j\to\infty}\int_{\R^n\times G(k,n)}f\,dT_j(x,\omega)=0.
\end{aligned}
\label{eq:3.19}
\end{equation}
\end{proof}

许多关于光滑极小曲面的早期结果, 都可以通过在 \eqref{eq:3.16} 中选取适当的向量场
而推广到驻定变分流形. 下面我们给出其中两个推广的论证: 坐标函数的调和性
与单调性公式. 这些结果应当与第 1 章的结果, 特别是\propref{prop:1.7}、1.12 与 1.15
作比较.

若对任意具有紧支集的光滑函数 $\eta$ 都有
\begin{equation}
\int_{\R^n\times G(k,n)}\langle\nabla_\omega\eta,\nabla_\omega u\rangle
 \,dT(x,\omega)=0,
\label{eq:3.20}
\end{equation}
则称 $\R^n$ 上的函数 $u$ 在变分流形 $T$ 上是\textbf{弱调和}的. 同样, 若
对任意具有紧支集的非负光滑函数 $\eta$ 都有
\begin{equation}
\int_{\R^n\times G(k,n)}\langle\nabla_\omega\eta,\nabla_\omega u\rangle
 \,dT(x,\omega)\leq0,
\label{eq:3.21}
\end{equation}
则称 $u$ 在 $T$ 上是\textbf{弱次调和}的.

\begin{proposition}[坐标函数的调和性]
\label{prop:3.6}
变分流形 $T^k\subset\R^n$ 是驻定的, 当且仅当坐标函数 $x_i$ 都是弱调和的.
\end{proposition}

\begin{proof}
设 $\eta$ 是具有紧支集的光滑函数, 令 $e_i=\nabla x_i$. 对任意单位向量
$E_j$ 有 $\nabla_{E_j}(\eta e_i)=(\nabla_{E_j}\eta)e_i$, 于是对任意
$k$ 维平面 $\omega$ 有
\begin{equation}
\dive_\omega(\eta e_i)=\langle\nabla_\omega\eta,e_i\rangle
 =\langle\nabla_\omega\eta,\nabla_\omega x_i\rangle,
\label{eq:3.22}
\end{equation}
其中 $\nabla_\omega\eta$ 是 $\nabla\eta$ 在 $\omega$ 上的投影. 因此
\begin{equation}
\int_{\R^n\times G(k,n)}\dive_\omega(\eta e_i)\,dT(x,\omega)
 =\int_{\R^n\times G(k,n)}
 \langle\nabla_\omega\eta,\nabla_\omega x_i\rangle\,dT(x,\omega).
\label{eq:3.23}
\end{equation}
由 \eqref{eq:3.23} 即得结论.
\end{proof}

同样, 单调性公式也可以推广到这个框架. 设 $\eta:\R\to\R$ 是非负函数, 记
$r=|x|$. 对任意向量 $E$, 有
\begin{equation}
\nabla_E x=E,
\label{eq:3.24}
\end{equation}
其中 $x=(x_1,\dots,x_n)$. 对以 $E_i$ 为标准正交基的 $k$ 维平面 $\omega$,
用 \eqref{eq:3.24} 计算得
\begin{equation}
\begin{aligned}
\dive_\omega(\eta(r)x)
&=\sum_{i=1}^k\langle E_i,\nabla_{E_i}(\eta(r)x)\rangle\\
&=k\eta(r)+\eta'(r)\langle\nabla_\omega r,x\rangle
 =k\eta(r)+r\eta'(r)|\nabla_\omega r|^2.
\end{aligned}
\label{eq:3.25}
\end{equation}
记 $\omega^N$ 为 $\omega$ 的正交 $(n-k)$ 维平面. 有
\begin{equation}
1=|\nabla r|^2=|\nabla_\omega r|^2+|\nabla_{\omega^N}r|^2.
\label{eq:3.26}
\end{equation}

\begin{proposition}[驻定变分流形的单调性]
\label{prop:3.7}
设 $T^k\subset\R^n$ 是驻定变分流形, $x_0\in\R^n$; 则对一切 $0<s<t$ 有
\begin{equation}
\begin{aligned}
&t^{-k}\mu_T(B_t(x_0))-s^{-k}\mu_T(B_s(x_0))\\
&\qquad=\int_{(B_t(x_0)\setminus B_s(x_0))\times G(k,n)}
 r^{-k}|\nabla_{\omega^N}r|^2\,dT(x,\omega).
\end{aligned}
\label{eq:3.27}
\end{equation}
\end{proposition}

\begin{proof}
经平移可设 $x_0=0$. 设 $\phi$ 是非负截断函数, $\phi'(s)\leq0$, 它在
$[0,\frac12]$ 上恒等于 $1$, 支集在 $[0,1]$ 中. 暂时固定 $s$, 令
$\eta(r)=\phi(\frac rs)$, 于是
\begin{equation}
r\eta'(r)=-s\frac{d}{ds}\Bigl(\phi\Bigl(\frac rs\Bigr)\Bigr).
\label{eq:3.28}
\end{equation}
由于 $T$ 驻定, 对 \eqref{eq:3.25} 积分并利用 \eqref{eq:3.26} 得
\begin{equation}
\begin{aligned}
0&=\int\dive_\omega(\eta(r)x)\,dT(x,\omega)\\
&=\int\bigl(k\eta(r)+r\eta'(r)\bigr)\,dT(x,\omega)
 -\int r\eta'(r)|\nabla_{\omega^N}r|^2\,dT(x,\omega).
\end{aligned}
\label{eq:3.29}
\end{equation}
把 \eqref{eq:3.28} 代入 \eqref{eq:3.29} 得
\begin{equation}
\begin{aligned}
&\int k\phi\Bigl(\frac rs\Bigr)-s\frac{d}{ds}
 \Bigl(\phi\Bigl(\frac rs\Bigr)\Bigr)\,dT(x,\omega)\\
&\quad=-s\int\frac{d}{ds}\Bigl(\phi\Bigl(\frac rs\Bigr)\Bigr)
 |\nabla_{\omega^N}r|^2\,dT(x,\omega).
\end{aligned}
\label{eq:3.30}
\end{equation}
两边乘以 $s^{-k-1}$, 可把 \eqref{eq:3.30} 改写为
\begin{equation}
\begin{aligned}
&\frac{d}{ds}\Bigl(s^{-k}\int\phi\Bigl(\frac rs\Bigr)\,dT(x,\omega)\Bigr)\\
&\quad=s^{-k}\frac{d}{ds}
 \Bigl(\int\phi\Bigl(\frac rs\Bigr)|\nabla_{\omega^N}r|^2
 \,dT(x,\omega)\Bigr).
\end{aligned}
\label{eq:3.31}
\end{equation}
让 $\phi$ 递增到 $[0,1]$ 的特征函数, 并对 \eqref{eq:3.31} 应用单调收敛定理, 得
\begin{equation}
\begin{aligned}
\frac{d}{ds}\bigl(s^{-k}\mu_T(B_s)\bigr)
&=s^{-k}\frac{d}{ds}
 \Bigl(\int_{B_s\times G(k,n)}|\nabla_{\omega^N}r|^2\,dT(x,\omega)\Bigr)\\
&=\frac{d}{ds}\Bigl(\int_{B_s\times G(k,n)}
 r^{-k}|\nabla_{\omega^N}r|^2\,dT(x,\omega)\Bigr).
\end{aligned}
\label{eq:3.32}
\end{equation}
等式 \eqref{eq:3.32} 在分布意义下成立, 也对几乎一切 $s$ 成立. 为看出这一点, 我们
利用 $\mu_T(B_s)$ 与 \eqref{eq:3.32} 右端积分的单调性. 把 \eqref{eq:3.32} 从 $s$ 积分到
$t$ 得
\begin{equation}
\begin{aligned}
&t^{-k}\mu_T(B_t)-s^{-k}\mu_T(B_s)\\
&\quad=\int_{(B_t\setminus B_s)\times G(k,n)}
 r^{-k}|\nabla_{\omega^N}r|^2\,dT(x,\omega).
\end{aligned}
\label{eq:3.33}
\end{equation}
\end{proof}

在第 1 章中, 我们为极小子流形定义了密度. 现在可以把它推广到驻定变分流形.
即: 给定驻定 $k$-变分流形 $T^k\subset\R^n$、$x_0\in\R^n$ 与 $s>0$,
定义密度为
\begin{equation}
\Theta_{x_0}(s)=\frac{\mu_T(B_s(x_0))}{\Vol(B_s\subset\R^k)}.
\label{eq:3.34}
\end{equation}
对\propref{prop:1.12} 中的论证作简单修改即可推出: $\Theta_{x_0}(s)$ 是 $s$ 的
单调非降函数.

同样, 更一般的平均值不等式也可以推广到驻定变分流形. 所需的修改留给读者.

\begin{proposition}[驻定变分流形的平均值不等式]
\label{prop:3.8}
若 $T\subset\R^n$ 是驻定 $k$-变分流形, $x_0\in\R^n$, $f$ 是 $T$ 上的
非负弱次调和函数, 且 $0<s<t$, 则
\begin{equation}
t^{-k}\int_{B_t(x_0)}f\,d\mu_T
 \geq s^{-k}\int_{B_s(x_0)}f\,d\mu_T.
\label{eq:3.35}
\end{equation}
\end{proposition}

至此我们讨论的都是一般变分流形的理论. 来自光滑子流形的变分流形具有
多得多的额外结构. 特别地, 相应的测度支撑在光滑子流形上, 而这些子流形在
每一点都有切空间. 称为可求长变分流形的那一类变分流形具有类似的额外结构.

记 $\mathcal{H}^k$ 为 $k$ 维 Hausdorff 测度.

\begin{definition}
\label{def:3.9}
若 $S\subset S_0\cup S_1$, 其中 $\mathcal{H}^k(S_0)=0$ 且 $S_1$ 是
$\R^k$ 在 Lipschitz 映射下的像, 则称集合 $S\subset\R^n$ 是
\textbf{$k$-可求长的}.

更一般地, 若 $S\subset\bigcup_{\ell\geq0}S_\ell$, 其中
$\mathcal{H}^k(S_0)=0$, 且对 $\ell\geq1$ 每个 $S_\ell$ 都是 $\R^k$ 在
Lipschitz 映射下的像, 则称 $S$ 是\textbf{可数 $k$-可求长的}.
\end{definition}

处理可求长集的一个好处是: 由 Rademacher 定理 \cite{Fe}, 它们在几乎每一点
都有切空间. 于是我们可以像对光滑子流形那样, 给可求长集配上一个变分流形.
可求长变分流形就是支撑在可求长集上的变分流形.

\begin{definition}[可求长变分流形]
\label{def:3.10}
设 $S$ 是 $\R^n$ 的可数 $k$-可求长子集, 且 $\mathcal{H}^k(S)<\infty$;
设 $\theta$ 是 $S$ 上取正值、且局部 $\mathcal{H}^k$ 可积的函数. 令 $T$ 等于
关联于集合 $S$ 的变分流形(完全如同 $S$ 是光滑子流形那样). 相应的变分流形
$T'=\theta T$ 称为\textbf{可求长变分流形}. 若 $\theta$ 取整数值, 则 $T'$
称为\newterm{integral-varifold}{整变分流形}{integral varifold}.
\end{definition}

此后, $\Sigma^k\subset\R^n$ 总表示驻定可求长 $k$-变分流形. 每个这样的
$\Sigma$ 都伴有一个 Radon 测度 $\mu_\Sigma$; 稍微滥用记号, 我们也用
$\Sigma$ 表示该测度所支撑的(Hausdorff $k$ 维可求长)集合. 我们作标准假设:
密度在支集上至少为 $1$. 注意这类广义极小子流形包含了配以内在黎曼度量的
嵌入极小子流形的情形.

\mtpara{变分流形距离}
固定一个闭流形 $M$, 令
\[
\Pi:G_kM\to M
\]
为(无向)$k$ 维平面的 Grassmann 丛, 即每条纤维 $\Pi^{-1}(p)$ 是 $M$ 在
$p$ 处切空间的一切 $k$ 维线性子空间构成的集合. 由于 $G_kM$ 紧, 我们可以
选取 $G_kM$ 上一切连续函数中关于上确界范数稠密的一个可数子集 $\{h_n\}$,
其中每个函数的上确界范数不超过 $1$.\footnote{这是 Stone--Weierstrass
定理的推论; 见 \cite{R} 第 213 页推论 35.}
若 $(X_0,F_0)$ 与 $(X_1,F_1)$ 是两个紧的(不必连通的)曲面 $X_0,X_1$,
并带有可测映射
\[
F_i:X_i\to G_kM,
\]
使得每个 $f_i=\Pi\circ F_i$ 属于 Sobolev 空间 $W^{1,2}(X_i,M)$(即平方
可积且梯度也平方可积的映射), 而 $J_{f_i}$ 是 $f_i$ 的 Jacobi 因子, 则
它们之间的\textbf{变分流形距离}按定义为
\begin{equation}
d_V(F_0,F_1)=\sum_n2^{-n}
 \Bigl|\int_{X_0}h_n\circ F_0\,J_{f_0}
 -\int_{X_1}h_n\circ F_1\,J_{f_1}\Bigr|.
\label{eq:3.36}
\end{equation}
由此容易看出: 面积一致有界的序列 $X_i=(X_i,F_i)$ 收敛于 $(X,F)$, 当且
仅当它弱收敛, 即对一切 $h\in C^0(G_2M)$ 有
\begin{equation}
\int_{X_i}h\circ F_i\,J_{f_i}\to\int_X h\circ F\,J_f.
\label{eq:3.37}
\end{equation}
例如, 当 $M$ 是三维流形时, $G_2M$、$G_1M$ 与 $T^1M/\{\pm v\}$ 同构
(这里 $T^1M$ 是单位切丛). 若 $\Sigma_i$ 是 $M$ 中一列闭浸入曲面, 在通常
的 $C^k$ 拓扑下收敛于闭曲面 $\Sigma$, 则我们可以把每个曲面嵌入
$T^1M/\{\pm v\}\equiv G_2M$ 中: 把每一点映为该曲面的单位法向量 $\pm N$.
容易看出, 带这些嵌入映射的曲面在变分流形距离下收敛. 更一般地, 若 $X$ 是
紧曲面, $f:X\to M$ 是 $W^{1,2}$ 映射(此时不再假定 $M$ 是三维的), 则令
$F:X\to G_2M$ 由``$F(x)$ 是线性子空间 $df(T_xX)$''给出(当 $M$ 是三维时,
可以把该映射的像看作落在 $T^1M/\{\pm v\}$ 中). 严格说来, 这只在 $J_f$
非零的可测集上有定义; 我们把它任意延拓到整个 $X$, 因为由
$h\mapsto\int_X h\circ F\,J_f$ 给出的 $G_2M$ 上的 Radon 测度与延拓方式
无关.

% ================================================================
\section{Sobolev 不等式}
\label{sec:3.2}

接下来我们证明 Michael 与 Simon \cite{MiSi} 的 Sobolev 不等式:

\begin{theorem}
\label{thm:3.11}
存在仅依赖于 $n$ 的常数 $c_n$, 使得若 $u\geq0$ 是超曲面
$\Sigma\subset\R^{n+1}$ 上具有紧支集的 Lipschitz 函数, 则
\begin{equation}
\Bigl(\int_\Sigma u^{\frac{n}{n-1}}\Bigr)^{\frac{n-1}{n}}
 \leq c_n\int_\Sigma\bigl(|\nabla u|+u|H|\bigr).
\label{eq:3.38}
\end{equation}
\end{theorem}

上述 Sobolev 不等式事实上是一族不等式的极端情形. 然而, 如下文所见,
所有这些不等式都可以由\thmref{thm:3.11} 与 Hölder 不等式推出.

\begin{corollary}
\label{cor:3.12}
存在仅依赖于 $n$ 的常数 $c_n$, 使得若 $u\geq0$ 是超曲面
$\Sigma\subset\R^{n+1}$ 上具有紧支集的 Lipschitz 函数, 且 $p\in[1,n)$,
则
\begin{equation}
\|u\|_{L^{\frac{np}{n-p}}}
 \leq c_n\Bigl(\frac{(n-1)p}{n-p}\|\nabla u\|_{L^p}
 +\|uH\|_{L^p}\Bigr).
\label{eq:3.39}
\end{equation}
\end{corollary}

\begin{proof}
给定 $r>1$, 对函数 $u^r$ 应用\thmref{thm:3.11} 得
\begin{equation}
\|u\|_{L^{\frac{nr}{n-1}}}^r
 =\Bigl(\int_\Sigma u^{\frac{nr}{n-1}}\Bigr)^{\frac{n-1}{n}}
 \leq c_n\int_\Sigma\bigl(ru^{r-1}|\nabla u|+u^r|H|\bigr).
\label{eq:3.40}
\end{equation}
以指数 $\frac{nr}{(n-1)(r-1)}$ 与 $\frac{nr}{n+r-1}$ 应用 Hölder 不等式得
\[
\begin{aligned}
\int_\Sigma u^{r-1}|\nabla u|
&\leq\Bigl(\int_\Sigma u^{\frac{nr}{n-1}}\Bigr)^{\frac{(n-1)(r-1)}{nr}}
 \Bigl(\int_\Sigma|\nabla u|^{\frac{nr}{n+r-1}}\Bigr)^{\frac{n+r-1}{nr}},\\
\int_\Sigma u^r|H|
&\leq\Bigl(\int_\Sigma u^{\frac{nr}{n-1}}\Bigr)^{\frac{(n-1)(r-1)}{nr}}
 \Bigl(\int_\Sigma(u|H|)^{\frac{nr}{n+r-1}}\Bigr)^{\frac{n+r-1}{nr}}.
\end{aligned}
\]
用这两个界估计 \eqref{eq:3.40} 右端的各项, 再除以
\[
\Bigl(\int_\Sigma u^{\frac{nr}{n-1}}\Bigr)^{\frac{(n-1)(r-1)}{nr}}
 =\|u\|_{L^{\frac{nr}{n-1}}}^{r-1},
\]
便得到
\[
\|u\|_{L^{\frac{nr}{n-1}}}
 \leq c_n\Bigl(r\|\nabla u\|_{L^{\frac{nr}{n+r-1}}}
 +\|uH\|_{L^{\frac{nr}{n+r-1}}}\Bigr).
\]
令 $p=\frac{nr}{n+r-1}$, 于是 $r=\frac{(n-1)p}{n-p}$ 且
$\frac{nr}{n-1}=\frac{np}{n-p}$, 即得推论.
\end{proof}

Michael--Simon Sobolev 不等式被 Hoffman 与 Spruck \cite{HoSp} 推广到了
一般黎曼流形的子流形.

\mtpara{Sobolev 不等式的证明}
我们需要下述覆盖引理(\cite{Si3} 中的定理 3.3):

\begin{lemma}
\label{lem:3.13}
若 $\mathcal{B}$ 是度量空间中一族闭球, 且
\[
\sup\{\diam(B)\mid B\in\mathcal{B}\}<\infty,
\]
则存在两两不相交的子族 $\mathcal{B}'\subset\mathcal{B}$, 使得
\[
\bigcup_{B\in\mathcal{B}}\subset\bigcup_{B\in\mathcal{B}'}5B,
\]
其中 $5B$ 是与 $B$ 同心、半径为 $B$ 的半径五倍的球.
\end{lemma}

\begin{proof}
令 $d=\sup\{\diam(B)\mid B\in\mathcal{B}\}$, 并把 $\mathcal{B}$ 分成互不
相交的族 $\mathcal{B}_1,\mathcal{B}_2,\dots$, 其中
\[
\mathcal{B}_j=\{B\in\mathcal{B}\mid
 \diam(B)\in(2^{-j}d,2^{1-j}d)\}.
\]
取 $\mathcal{B}_1$ 中任一两两不相交的极大子族 $\mathcal{B}_1'$. 其次, 取
\[
\{B\in\mathcal{B}_2\mid \text{对每个 }B'\in\mathcal{B}_1'
 \text{ 都有 }B\cap B'=\emptyset\}
\]
中任一两两不相交的极大子族 $\mathcal{B}_2'$. 如此归纳地定义: 取
\[
\{B\in\mathcal{B}_j\mid \text{对每个 }
B'\in\bigcup_{i<j}\mathcal{B}_i'\text{ 都有 }B\cap B'=\emptyset\}
\]
中任一两两不相交的极大子族 $\mathcal{B}_j'$. 现令
$\mathcal{B}'=\bigcup_i\mathcal{B}_i'$. 它显然是两两不相交的.

给定 $B\in\mathcal{B}$, 它必属于某个 $\mathcal{B}_j$. 若
$B\in\mathcal{B}_j'$, 则显然已经证完; 故可设 $B\notin\mathcal{B}_j'$. 由
族 $\mathcal{B}_j'$ 的极大性, 必存在 $i\leq j$ 与球
$B'\in\mathcal{B}_i$, 使得
\[
B\cap B'\neq\emptyset.
\]
由 $i\leq j$ 可得 $\diam B\leq2\diam B'$. 于是由三角不等式,
$B\subset5B'$, 证毕.
\end{proof}

我们将遵循 Leon Simon 在 \cite{Si3} 中对 Sobolev 不等式的证明. 设
$\Sigma\subset\R^{n+1}$ 是超曲面, 定义 $J(r)$ 为
\begin{equation}
J(r)=\frac{1}{\omega_nr^n}\int_{B_r\cap\Sigma}h,
\label{eq:3.41}
\end{equation}
其中 $\omega_n$ 是 $\R^n$ 中单位球的体积, $h$ 是具有紧支集的非负函数.
\lemref{lem:1.18} 的平均值不等式给出
\begin{equation}
r^{n+1}J'(r)\geq\frac{1}{\omega_n}\int_{B_r\cap\Sigma}
 \langle x,\nabla h+hHN\rangle
 \geq-\frac{r}{\omega_n}\int_{B_r\cap\Sigma}
 \bigl(|\nabla h|+h|H|\bigr).
\label{eq:3.42}
\end{equation}
对此积分, 得: 对 $0<\sigma<\rho$ 有
\begin{equation}
J(\sigma)\leq J(\rho)+\int_\sigma^\rho
 \Bigl(r^{-n}\frac{1}{\omega_n}\int_{B_r\cap\Sigma}
 \bigl(|\nabla h|+h|H|\bigr)\Bigr)dr.
\label{eq:3.43}
\end{equation}

下一个要素是一个微积分引理(参看 \cite{Si3} 中的 18.7):

\begin{lemma}
\label{lem:3.14}
设 $f,g\geq0$ 有界、单调非降, 且
\begin{equation}
1\leq\limsup_{\sigma\to0}\frac{f(\sigma)}{\sigma^n},
\label{eq:3.44}
\end{equation}
并且对任意 $0<\sigma<\rho$ 有
\begin{equation}
\frac{f(\sigma)}{\sigma^n}\leq\frac{f(\rho)}{\rho^n}
 +\int_0^\rho t^{-n}g(t)\,dt.
\label{eq:3.45}
\end{equation}
令 $f_\infty=\lim_{t\to\infty}f(t)$, $\rho_0=2(f_\infty)^{1/n}$. 则存在
$\rho\in(0,\rho_0)$, 使得
\begin{equation}
f(5\rho)\leq\frac{5^n}{2}\rho_0\,g(\rho).
\label{eq:3.46}
\end{equation}
\end{lemma}

\begin{proof}
用反证法, 设对每个 $\rho\in(0,\rho_0)$ 都有 \eqref{eq:3.46} 不成立. 把它与
\eqref{eq:3.45} 结合得
\begin{equation}
\begin{aligned}
\sup_{\sigma\in(0,\rho_0)}\frac{f(\sigma)}{\sigma^n}
&\leq\frac{f(\rho_0)}{\rho_0^n}+\int_0^{\rho_0}t^{-n}g(t)\,dt\\
&\leq\frac{f_\infty}{\rho_0^n}
 +\frac{2}{5^n\rho_0}\int_0^{\rho_0}t^{-n}f(5t)\,dt.
\end{aligned}
\label{eq:3.47}
\end{equation}
为估计第二个积分, 作变量替换 $s=5t$, 得
\[
\begin{aligned}
\frac{2}{5^n\rho_0}\int_0^{\rho_0}t^{-n}f(5t)\,dt
&=\frac{2}{5\rho_0}\int_0^{5\rho_0}s^{-n}f(s)\,ds\\
&=\frac{2}{5\rho_0}\int_0^{\rho_0}s^{-n}f(s)\,ds
 +\frac{2}{5\rho_0}\int_{\rho_0}^{5\rho_0}s^{-n}f(s)\,ds\\
&\leq\frac25\sup_{\sigma\in(0,\rho_0)}\frac{f(\sigma)}{\sigma^n}
 +\frac{2}{5(n-1)}\rho_0^{-n}f_\infty.
\end{aligned}
\]
把它代回 \eqref{eq:3.47} 得
\[
\sup_{\sigma\in(0,\rho_0)}\frac{f(\sigma)}{\sigma^n}
 \leq\frac{f_\infty}{\rho_0^n}
 +\frac25\sup_{\sigma\in(0,\rho_0)}\frac{f(\sigma)}{\sigma^n}
 +\frac{2}{5(n-1)}\frac{f_\infty}{\rho_0^n},
\]
两边减去右端第二项得
\begin{equation}
\frac35\sup_{\sigma\in(0,\rho_0)}\frac{f(\sigma)}{\sigma^n}
 \leq\frac{5n-3}{5(n-1)}\frac{f_\infty}{\rho_0^n}
 =\frac{5n-3}{5(n-1)}2^{-n}.
\label{eq:3.48}
\end{equation}
最后一个等号用了 $\rho_0$ 的定义. 乘以 $5/3$, 并注意当 $n\geq2$ 时
$\frac{5n-3}{5(n-1)}\leq7/5$, 得
\begin{equation}
\sup_{\sigma\in(0,\rho_0)}\frac{f(\sigma)}{\sigma^n}
 \leq\frac73\,2^{-n}\leq\frac7{12}<1,
\label{eq:3.49}
\end{equation}
但这与 \eqref{eq:3.44} 矛盾, 从而引理得证.
\end{proof}

若 $y\in\Sigma$ 是满足 $h(y)\geq1$ 的点, 则取
\begin{gather*}
f(\rho)=\frac{1}{\omega_n}\int_{B_\rho(y)\cap\Sigma}h,\\
g(\rho)=\frac{1}{\omega_n}\int_{B_\rho(y)\cap\Sigma}
 \bigl(|\nabla h|+h|H|\bigr)
\end{gather*}
时, \eqref{eq:3.44} 与 \eqref{eq:3.45} 都成立. (我们用 \eqref{eq:3.43} 说明 \eqref{eq:3.45} 成立; 至于
\eqref{eq:3.44}, 注意 $\lim_{\rho\to0}f(\rho)$ 就是 $h(y)$ 乘以 $\Sigma$ 在 $y$
处的重数.) 如此选取 $f$, 有
$f_\infty=\frac{1}{\omega_n}\int_\Sigma h<\infty$(因为 $h$ 具有紧支集),
于是
\[
\rho_0=2f_\infty^{1/n}
 =2\Bigl(\frac{1}{\omega_n}\int_\Sigma h\Bigr)^{1/n}.
\]
因此\lemref{lem:3.14} 给出半径 $\rho_y$, 使得
\begin{equation}
\int_{B_{5\rho_y}(y)\cap\Sigma}h
 \leq\frac{5^n}{\omega_n^{1/n}}
 \Bigl(\int_\Sigma h\Bigr)^{1/n}
 \int_{B_{\rho_y}(y)\cap\Sigma}\bigl(|\nabla h|+h|H|\bigr).
\label{eq:3.50}
\end{equation}
设 $S\subset\Sigma$ 是 $h\geq1$ 的集合:
\[
S=\{y\in\Sigma\mid h(y)\geq1\}.
\]
利用五倍覆盖引理(\lemref{lem:3.13}), 我们可以选取互不相交的球 $B^j$, 使之满足
\eqref{eq:3.50}, 并且
\begin{equation}
S\subset\bigcup_j5(B^j).
\label{eq:3.51}
\end{equation}
于是
\begin{equation}
\int_S h\leq\frac{5^n}{\omega_n^{1/n}}
 \Bigl(\int_\Sigma h\Bigr)^{1/n}
 \int_\Sigma\bigl(|\nabla h|+h|H|\bigr).
\label{eq:3.52}
\end{equation}

固定 $\epsilon>0$ 以及单调函数 $\gamma(t)$, 使得 $t\geq\epsilon$ 时
$\gamma(t)=1$, $t\leq0$ 时 $\gamma(t)=0$. 对任意 $t>0$, 令
\begin{equation}
A_t=\{x\in\Sigma\mid u(x)>t\}.
\label{eq:3.53}
\end{equation}
设 $\mu$ 是 $\Sigma$ 上的体积测度. 对 $h(x)=\gamma(u(x)-t)$ 应用 \eqref{eq:3.52} 得
\begin{equation}
\mu(A_{t+\epsilon})\leq\frac{5^n}{\omega_n^{1/n}}
 \bigl(\mu(A_t)\bigr)^{1/n}
 \int_\Sigma\bigl(|\nabla u|\gamma'(u-t)+\gamma(u-t)|H|\bigr).
\label{eq:3.54}
\end{equation}
两边乘以 $(t+\epsilon)^{\frac{1}{n-1}}$ 得
\begin{equation}
\begin{aligned}
&(t+\epsilon)^{\frac{1}{n-1}}\mu(A_{t+\epsilon})\\
&\quad\leq\frac{5^n}{\omega_n^{1/n}}
 \Bigl((t+\epsilon)^{\frac{n}{n-1}}\mu(A_t)\Bigr)^{1/n}
 \int_\Sigma\bigl(|\nabla u|\gamma'(u-t)+\gamma(u-t)|H|\bigr)\\
&\quad\leq\frac{5^n}{\omega_n^{1/n}}
 \Bigl(\int_\Sigma(u+\epsilon)^{\frac{n}{n-1}}\Bigr)^{1/n}
 \Bigl(\int_\Sigma|\nabla u|\gamma'(u-t)+\int_{A_t}|H|\Bigr),
\end{aligned}
\label{eq:3.55}
\end{equation}
其中最后一个不等式用了对任意非负函数 $v$ 都成立的显然不等式
$t\mu(\{v>t\})\leq\int v$.

我们还需要最后一个引理, 它是 Fubini 定理的标准应用:

\begin{lemma}
\label{lem:3.15}
设 $\nu$ 是空间 $X$ 上的测度, $f\geq0$ 属于 $L^1(\nu)$,
$A_t=\{x\mid f(x)>t\}$, $\alpha>0$. 则
\begin{equation}
\frac1\alpha\int f^\alpha\,d\nu
 =\int_0^\infty t^{\alpha-1}\nu(A_t)\,dt.
\label{eq:3.56}
\end{equation}
由此可知, 若 $t_0\geq0$, 则
\begin{equation}
\frac1\alpha\int_{A_{t_0}}(f^\alpha-t_0^\alpha)\,d\nu
 =\int_{t_0}^\infty t^{\alpha-1}\nu(A_t)\,dt.
\label{eq:3.57}
\end{equation}
\end{lemma}

\begin{proof}
该引理由对乘积空间 $X\times[t_0,\infty)$ 应用 Fubini 定理得到.
\end{proof}

现在可以完成 Sobolev 不等式的证明.

\begin{proof}[\thmref{thm:3.11} 的证明]
对 \eqref{eq:3.55} 关于 $t$ 积分, 并对第一项应用\lemref{lem:3.15}(取 $\nu=\mu$,
$t_0=\epsilon$, $\alpha-1=\frac{1}{n-1}$, 即 $\alpha=\frac{n}{n-1}$), 得
\begin{equation}
\begin{aligned}
&\frac{n-1}{n}\int_{A_\epsilon}
 \bigl(u^{\frac{n}{n-1}}-\epsilon^{\frac{n}{n-1}}\bigr)\\
&\quad\leq\frac{5^n}{\omega_n^{1/n}}
 \Bigl(\int_\Sigma(u+\epsilon)^{\frac{n}{n-1}}\Bigr)^{\frac1n}
 \int_0^\infty\Bigl(\int_\Sigma|\nabla u|\gamma'(u-t)
 +\int_{A_t}|H|\Bigr)dt\\
&\quad=\frac{5^n}{\omega_n^{1/n}}
 \Bigl(\int_\Sigma(u+\epsilon)^{\frac{n}{n-1}}\Bigr)^{\frac1n}
 \int_\Sigma\bigl(|\nabla u|+u|H|\bigr),
\end{aligned}
\label{eq:3.58}
\end{equation}
其中等号用了微积分基本定理(以及 Fubini 定理交换积分次序)来计算
$\gamma'$ 项, 并用\lemref{lem:3.15}(取 $\nu=|H|\mu$, $t_0=0$, $\alpha=1$)计算
最后一项. 令 $\epsilon\to0$, 得
\begin{equation}
\frac{n-1}{n}\int u^{\frac{n}{n-1}}
 \leq\frac{5^n}{\omega_n^{1/n}}
 \Bigl(\int_\Sigma u^{\frac{n}{n-1}}\Bigr)^{1/n}
 \int_\Sigma\bigl(|\nabla u|+u|H|\bigr),
\label{eq:3.59}
\end{equation}
这就给出定理, 其中
\[
c_n=\frac{n}{n-1}\frac{5^n}{\omega_n^{1/n}}.
\]
\end{proof}

% ================================================================
\section{弱 Bernstein 型定理}
\label{sec:3.3}

我们将下述结果称为\textbf{弱 Bernstein 型定理}.

\begin{theorem}[Colding--Minicozzi, \cite{CM16}]
\label{thm:3.16}
若 $\Sigma^k\subset\R^n$ 是 $k$-可求长驻定变分流形, 其(体积)密度几乎
处处至少为 $1$, 且以 $\mathrm{V}_\Sigma$ 为上界, 则 $\Sigma$ 必含于某个
维数不超过
$(k+1)\frac{k}{k-1}\mathrm{e}^8\,2^{k+4}\mathrm{V}_\Sigma$ 的仿射子空间
之中.
\end{theorem}

\thmref{thm:3.16} 中的常数不是最优的, 也不是从 \cite{CM16} 的论证所能得到的最佳
常数. 然而, 几何依赖性是最优的(即界关于 $\mathrm{V}_\Sigma$ 是线性的,
这一点最优).

注意\thmref{thm:3.16} 用 $\Sigma$ 的体积给出了其上线性无关坐标函数个数的一个界.
当 $\mathrm{V}_\Sigma$ 充分小时, Allard 定理表明 $\Sigma^k$ 是平坦的,
因而只有 $k$ 个独立的坐标函数. 我们指出, 由\thmref{thm:3.16} 可知 Allard 正则性
定理中的 $\delta$ 与 $n$ 无关.

一个校准论证(见 \cite{Fe} 或 \cite{Mr}), 参看 \eqref{eq:1.15} 与 \eqref{eq:1.16}, 表明
$\C^n=\R^{2n}$ 的复子流形是绝对面积最小化的(从而是极小的). 利用这一点
容易看出, 仿射代数簇就是具有有界密度的驻定可求长变分流形的例子.

前面我们看到, 坐标函数在 $\Sigma$ 上是弱调和的(\propref{prop:3.6}; 参看
\propref{prop:1.7}). \thmref{thm:3.16} 将由 $\Sigma$ 上某类调和函数空间维数的界推出, 这类
函数就是所谓的多项式增长的调和函数.

\begin{definition}[多项式增长的调和函数]
\label{def:3.17}
我们定义 $\mathcal{H}_d(\Sigma)$ 为增长阶不超过 $d$ 的多项式增长的调和
函数构成的线性空间. 即: $u\in\mathcal{H}_d(\Sigma)$, 若 $u$ 调和且存在
$C<\infty$, 使得 $|u(x)|\leq C(1+|x|^d)$.
\end{definition}

按此定义, 坐标函数 $x_i$ 属于 $\mathcal{H}_1(\Sigma)$. 例如, 在以
$x_3$ 轴为中心的悬链面 $\Sigma_c^2\subset\R^3$ 上, 函数 $x_3$ 的增长比
内蕴距离的任何幂次都慢; 然而对任何 $d<1$, 它不属于
$\mathcal{H}_d(\Sigma_c^2,\R)$.

\begin{theorem}[Colding--Minicozzi, \cite{CM16}]
\label{thm:3.18}
设 $\Sigma^k$ 是驻定 $k$-可求长变分流形, 其密度几乎处处至少为 $1$, 且以
$\mathrm{V}_\Sigma<\infty$ 为上界. 则对任意 $d\geq1$,
\begin{equation}
\dim\mathcal{H}_d(\Sigma)\leq C\,\mathrm{V}_\Sigma\,d^{k-1},
\qquad\text{其中 } C=(k+1)\frac{k}{k-1}\mathrm{e}^8\,2^{k+4}.
\label{eq:3.60}
\end{equation}
\end{theorem}

这里有限维数结果的证明由以下两个独立的步骤组成. 首先, 我们把问题化归为
估计 $\Sigma$ 上 $L^2(B_r\cap\Sigma)$-标准正交的调和函数的个数, 这些函数
在 $L^2(B_{\Omega r}\cap\Sigma)$ 范数下的界对一切 $r>0$ 与
$\Omega>1$ 一致. 其次, 我们利用调和函数空间的线性性构造一个``Bergman 核''. 即: 若
$u_1,\dots,u_{\mathcal{N}}$ 是这些调和函数, 令
\begin{equation}
K(x)=\sum_{i=1}^{\mathcal{N}}|u_i(x)|^2.
\label{eq:3.61}
\end{equation}
把一些标准的线性代数与平均值不等式结合起来, 我们得到 $K(x)$ 的逐点界,
它依赖于 $\Omega$ 与 $L^2(B_{\Omega r}\cap\Sigma)$ 范数, 但与
$\mathcal{N}$ 无关. 把它在 $B_r\cap\Sigma$ 上积分, 就得到 $\mathcal{N}$
的一个依赖于 $\Omega$ 的界.

为了得到 $\mathcal{H}_d(\Sigma)$ 的维数关于 $d$ 的多项式界, 我们取
$\Omega$ 约为 $1+\frac1d$. 所得的估计关于 $d$ 是多项式的, 其次数既依赖
于平均值不等式在边界处的爆破速率, 也依赖于体积测度的正则性.

关于这一方向的有关结果与更一般的结果, 见 \cite{CM16}. 关于其他框架下
多项式增长的调和函数的结果, 见 \cite{CM14}、\cite{CM15} 及其中的参考文献.
最后, 关于向平均曲率流的推广, 见 Calle \cite{Ca}.

% ================================================================
\section{一般构造}
\label{sec:3.4}

本节中我们遵循 \cite{CM16} 研究函数空间的增长性质与独立性. 首先, 下面的
引理说明关于函数增长率的假设是有用的.

\begin{lemma}
\label{lem:3.19}
设 $f_1,\dots,f_{2s}$ 是 $(0,\infty)$ 上的非负单调非降函数, 其中没有
一个恒为零, 且存在 $d_0,K>0$, 使得对一切 $i$ 有
\begin{equation}
f_i(r)\leq K(r^{d_0}+1).
\label{eq:3.62}
\end{equation}
则对任意 $\Omega>1$, 存在其中 $s$ 个函数
$f_{\alpha_1},\dots,f_{\alpha_s}$ 以及无穷多个整数 $m\geq1$, 使得对
$i=1,\dots,s$ 有
\begin{equation}
f_{\alpha_i}(\Omega^{m+1})\leq\Omega^{2d_0}f_{\alpha_i}(\Omega^m).
\label{eq:3.63}
\end{equation}
\end{lemma}

\begin{proof}
由于这些函数单调非降且都不恒为零, 我们可以假定存在 $R>0$, 使得对任意
$r>R$ 与一切 $i$ 都有 $f_i(r)>0$.

我们将证明: 存在无穷多个 $m$, 使得 $\{f_i\}$ 中存在一个 $s$ 元子集(该
子集可以随 $m$ 变化)满足 \eqref{eq:3.63}. 这就足以证明引理; 因为 $2s$ 个函数的
$s$ 元子集只有有限多个, 其中必有一个被无穷多次重复.

对 $r>R$, 注意
\begin{equation}
g(r)=\prod_{i=1}^{2s}f_i(r)\leq K^{2s}(r^{d_0}+1)^{2s},
\label{eq:3.64}
\end{equation}
且 $g$ 是正的非降函数. 假设只有有限多个
$m\geq\frac{\log R}{\log\Omega}$ 满足 \eqref{eq:3.63}. 设 $m_0-1$ 是其中最大的
$m$; 则对一切 $j\geq1$ 有
\begin{equation}
\Omega^{2d_0(s+1)}g(\Omega^{m_0+j-1})<g(\Omega^{m_0+j}).
\label{eq:3.65}
\end{equation}
迭代此式并应用 \eqref{eq:3.64} 得: 对任意 $j\geq1$ 有
\begin{equation}
\Omega^{2d_0(s+1)j}g(\Omega^{m_0})<g(\Omega^{m_0+j})
 \leq\tilde{c}\,(\Omega^j)^{2sd_0},
\label{eq:3.66}
\end{equation}
其中 $\tilde{c}=\tilde{c}(s,m_0,\Omega,K)$. 由于 $\Omega>1$, 取 $j$
充分大即得矛盾.
\end{proof}

对 $r>0$, $B_r=\{|x|<r\}\subset\R^n$ 以及函数 $u$ 与 $v$, 令
\begin{equation}
I_u(r)=\int_{B_r\cap\Sigma}u^2
\label{eq:3.67}
\end{equation}
以及
\begin{equation}
J_r(u,v)=\int_{B_r\cap\Sigma}uv.
\label{eq:3.68}
\end{equation}
注意 $J_r$ 是内积, $I(r)$ 是相应的二次型. 此外, 若
$u\in\mathcal{H}_d(\Sigma)$, 则 $I_u(r)\leq C(1+r^{2d+k})$.

给定 $\mathcal{H}_d(\Sigma)$ 中一组线性无关的函数, 我们将构造单变量函数,
以反映这组函数的增长性质与独立性.

我们从两个定义开始. 第一个定义构造将要研究其增长性质的函数.

\begin{definition}[$w_{i,r}$ 与 $f_i$]
\label{def:3.20}
设 $u_1,\dots,u_s$ 是线性无关的函数. 对每个 $r>0$, 我们现在定义
$J_r$-正交的生成集 $w_{i,r}$ 以及函数 $f_i$. 令
$w_{1,r}=w_1=u_1$, $f_1(r)=I_{w_1}(r)$. 定义 $w_{i,r}$ 如下: 要求它
关于内积 $J_r$ 与 $u_j|B_r$($j<i$)正交, 并且
\begin{equation}
u_i=\sum_{j=1}^{i-1}\lambda_{ji}(r)u_j+w_{i,r}.
\label{eq:3.69}
\end{equation}
\end{definition}

注意, 若 $u_i|B_r$ 线性相关, 则 $\lambda_{ij}(r)$ 不是唯一确定的. 然而,
由于 $u_i$ 在 $\Sigma$ 上线性无关, 当 $r$ 充分大时 $\lambda_{ij}(r)$ 是
唯一确定的.

无论如何, 下述量对一切 $r>0$ 都有定义(事实上当 $r$ 充分大时为正)
\begin{equation}
f_i(r)=\int_{B_r}|w_{i,r}|^2.
\label{eq:3.70}
\end{equation}

在下一个命题中, 我们记录\defref{def:3.20} 中函数 $f_i$ 的一些关键性质.

\begin{proposition}[$f_i$ 的性质, \cite{CM16}]
\label{prop:3.21}
若 $u_1,\dots,u_s\in\mathcal{H}_d(\Sigma)$ 线性无关, 则\defref{def:3.20} 中的
$f_i$ 具有下述四条性质: 存在常数 $K>0$(依赖于集合 $\{u_i\}$), 使得对
$i=1,\dots,s$ 有
\begin{equation}
f_i(r)\leq K\bigl(r^{2d+k}+1\bigr),
\label{eq:3.71}
\end{equation}
\begin{equation}
\text{$f_i$ 是单调非降函数,}
\label{eq:3.72}
\end{equation}
\begin{equation}
\text{$f_i$ 非负, 且当 $r$ 充分大时为正,}
\label{eq:3.73}
\end{equation}
以及
\begin{equation}
f_i(r)=I_{w_{i,r}}(r),\qquad
\text{且当 } t<r \text{ 时 } f_i(t)\leq I_{w_{i,r}}(t).
\label{eq:3.74}
\end{equation}
\end{proposition}

\begin{proof}
首先注意 $f_i(r)\leq I_{u_i}(r)$; 于是得到 \eqref{eq:3.71}. 此外, 对 $s<r$ 有
\begin{equation}
\begin{aligned}
f_i(s)&=\int_{B_s}\Bigl|u_i-\sum_{j=1}^{i-1}\lambda_{ji}(s)u_j\Bigr|^2
 =I_{w_{i,s}}(s)\\
&\leq\int_{B_s}\Bigl|u_i-\sum_{j=1}^{i-1}\lambda_{ji}(r)u_j\Bigr|^2
 =I_{w_{i,r}}(s)\\
&\leq\int_{B_r}\Bigl|u_i-\sum_{j=1}^{i-1}\lambda_{ji}(r)u_j\Bigr|^2
 =I_{w_{i,r}}(r)=f_i(r),
\end{aligned}
\label{eq:3.75}
\end{equation}
其中 \eqref{eq:3.75} 的第一个不等式由 $w_{i,r}$ 对 $j<i$ 的 $u_j$ 的正交性得到,
\eqref{eq:3.75} 的第二个不等式由 $I$ 的单调性得到. 由 \eqref{eq:3.75} 与 $u_i$ 的线性
无关性, 我们得到 \eqref{eq:3.72} 与 \eqref{eq:3.73}. 最后, \eqref{eq:3.75} 也包含了 \eqref{eq:3.74}.
\end{proof}

在下一个命题中, 我们将对\defref{def:3.20} 中的函数 $f_i$ 应用\lemref{lem:3.19}:

\begin{proposition}[Colding--Minicozzi, \cite{CM16}]
\label{prop:3.22}
设 $u_1,\dots,u_{2s}\in\mathcal{H}_d(M)$ 线性无关. 给定 $\Omega>1$ 与
$m_0>0$, 存在 $m\geq m_0$、整数
$\ell\geq\frac12\Omega^{-4d-2k}s$, 以及 $u_i$ 张成的线性空间中的函数
$v_1,\dots,v_\ell$, 使得对 $i,j=1,\dots,\ell$ 有
\begin{equation}
J_{\Omega^{m+1}}(v_i,v_j)=\delta_{i,j}
\label{eq:3.76}
\end{equation}
并且
\begin{equation}
\frac12\Omega^{-4d-2k}\leq I_{v_i}(\Omega^m).
\label{eq:3.77}
\end{equation}
\end{proposition}

\begin{proof}
由\propref{prop:3.21} 的 \eqref{eq:3.71}、\eqref{eq:3.72} 与 \eqref{eq:3.74}, 对\defref{def:3.20} 中的 $f_i$ 应用
\lemref{lem:3.19}, 可知存在 $m\geq m_0$ 以及子集
$f_{\alpha_1},\dots,f_{\alpha_s}$, 使得对 $i=1,\dots,s$ 有
\begin{equation}
0<f_{\alpha_i}(\Omega^{m+1})
 \leq\Omega^{4d+2k}f_{\alpha_i}(\Omega^m).
\label{eq:3.78}
\end{equation}
设 $w_{\alpha_i,\Omega^{m+1}}$, $i=1,\dots,s$, 是\defref{def:3.20} 中相应的
$u_i$ 的线性组合.

考虑由函数 $w_{\alpha_i,\Omega^{m+1}}$ 张成的 $s$ 维线性空间, 其内积为
$J_{\Omega^{m+1}}$. 在这个空间上还有半正定的双线性形式
$J_{\Omega^m}$. 设 $v_1,\dots,v_s$ 是 $J_{\Omega^{m+1}}$ 的标准正交基,
且使 $J_{\Omega^m}$ 对角化. 我们现在就这两组基来计算 $J_{\Omega^m}$ 的
迹. 首先, 关于正交基 $w_{\alpha_i,\Omega^{m+1}}$, 由 \eqref{eq:3.74} 与 \eqref{eq:3.78} 得
\begin{equation}
s\Omega^{-4d-2k}\leq\sum_{i=1}^s
 \frac{I_{w_{\alpha_i,\Omega^{m+1}}}(\Omega^m)}
 {I_{w_{\alpha_i,\Omega^{m+1}}}(\Omega^{m+1})}.
\label{eq:3.79}
\end{equation}
由于迹与基的选取无关, 在标准正交基 $v_i$ 上计算便得
\begin{equation}
s\Omega^{-4d-2k}\leq\sum_{i=1}^sI_{v_i}(\Omega^m).
\label{eq:3.80}
\end{equation}
把它与
\begin{equation}
0\leq I_{v_i}(\Omega^m)\leq1
\label{eq:3.81}
\end{equation}
(后者由 $I$ 的单调性得到)结合, 可知至少有
$\ell\geq\frac s2\Omega^{-4d-2k}$ 个 $v_i$ 使得对其中的每一个都有
\begin{equation}
\frac12\Omega^{-4d-2k}\leq I_{v_i}(\Omega^m)
 \leq I_{v_i}(\Omega^{m+1})=1.
\label{eq:3.82}
\end{equation}
这就证明了命题.
\end{proof}

% ================================================================
\section{有限维数性}
\label{sec:3.5}

本节中, 我们将遵循 \cite{CM16} 说明如何估计满足平均值不等式的多项式增长
函数空间的维数. 维数的界关于增长速率是多项式的, 其指数由平均值不等式的
边界爆破决定.

沿用 \cite{CM16}, 若 $0<\epsilon\leq1$, $1\leq C_W<\infty$, 并且对任意
$0<\delta\leq\frac12$ 都存在 $R_0>0$, 使得对一切 $r\geq R_0$ 有
\begin{equation}
\Vol\bigl(B_r\setminus B_{(1-\delta)r}\cap\Sigma\bigr)
 \leq C_W\delta^\epsilon\Vol(B_r\cap\Sigma),
\label{eq:3.83}
\end{equation}
则称 $\Sigma$ 具有 \textbf{$\epsilon$-体积正则性}.

这有时也称为微加倍(micro-doubling).

例如, 由此立即看出任何 $k$ 维锥都具有 $1$-体积正则性. 我们这里不会用到,
但注意 \cite{CM16} 中的引理 3.3 表明具有体积加倍性质的测地度量空间满足
这种体积正则性; 有关结果与应用见 \cite{Ch}、\cite{NrPy}、\cite{NrTa}
与 \cite{Te}.

密度有上下界的驻定 $k$-可求长变分流形具有 $1$-体积正则性. 为看出这一点,
设 $\Sigma$ 是驻定 $k$-可求长变分流形, 其密度几乎处处至少为 $1$, 且
\begin{equation}
\mathrm{V}_\Sigma\equiv\lim_{r\to\infty}\Theta_0(r)<\infty.
\label{eq:3.84}
\end{equation}
给定 $0<\delta\leq\frac12$, 选取 $R_0$ 使得对 $R\geq R_0/2$ 有
\begin{equation}
\mathrm{V}_\Sigma-\Theta_0(R)<\delta\mathrm{V}_\Sigma,
\label{eq:3.85}
\end{equation}
于是对 $R\geq R_0$ 有
\begin{equation}
\begin{aligned}
\Vol\bigl(B_{(1-\delta)R}\cap\Sigma\bigr)
&=\Theta_0\bigl((1-\delta)R\bigr)\Vol(B_1\subset\R^k)(1-\delta)^kR^k\\
&\geq\mathrm{V}_\Sigma\Vol(B_1\subset\R^k)(1-\delta)^{k+1}R^k.
\end{aligned}
\label{eq:3.86}
\end{equation}
因此, 对 $R\geq R_0$ 有
\begin{equation}
\begin{aligned}
&\Vol\bigl(B_R\setminus B_{(1-\delta)R}\cap\Sigma\bigr)\\
&\qquad\leq\mathrm{V}_\Sigma\Vol(B_1\subset\R^k)
 \bigl(1-(1-\delta)^{k+1}\bigr)R^k\\
&\qquad\leq2(k+1)\delta\Vol(B_R\cap\Sigma).
\end{aligned}
\label{eq:3.87}
\end{equation}

下述命题将与\secref{sec:3.4}的结果结合使用:

\begin{proposition}[Colding--Minicozzi, \cite{CM16}]
\label{prop:3.23}
设 $\Sigma^k$ 是驻定 $k$-可求长变分流形, 其密度几乎处处至少为 $1$, 且以
$\mathrm{V}_\Sigma<\infty$ 为上界. 设 $0<a<1$ 固定, $r>2R_0$, 且
$v_1,\dots,v_{\mathcal{N}}$ 是调和的、$J_r$-标准正交的函数. 若对某个
$d\geq1$, 使得对任意 $R\geq R_0$, \eqref{eq:3.87} 对任意
$\delta\geq\frac{1}{4d}$ 成立, 并且对一切 $i$ 有
\begin{equation}
a\leq I_{v_i}\bigl((1-(2d)^{-1})r\bigr),
\label{eq:3.88}
\end{equation}
则
\begin{equation}
\mathcal{N}\leq C\,d^{k-1},
\label{eq:3.89}
\end{equation}
其中 $C=\mathrm{V}_\Sigma(k+1)\frac{k}{k-1}a^{-1}2^{k+1}$.
\end{proposition}

\begin{proof}
由于 $d\geq1$, 我们可以选取正整数 $N$, 使得 $d\leq N\leq2d$. 对每个
$x\in B_r$, 令
\begin{equation}
K(x)=\sum_{i=1}^{\mathcal{N}}|v_i|^2(x).
\label{eq:3.90}
\end{equation}
注意, 由于每个 $v_i\in L^2_{loc}(M)$, 由平均值不等式可知 $K(x)$ 必有限.
按构造, $K(x)$ 是对称双线性形式
\begin{equation}
(v,w)\mapsto\langle v,w\rangle(x)
\label{eq:3.91}
\end{equation}
的迹, 其中 $v,w$ 取遍 $v_i$ 张成的空间.

回忆对称矩阵总可以用正交基变换对角化. 因此, 给定 $x\in B_r$, 我们可以
为内积空间 $(\Span\{v_i\},J_r)$ 选取一组新的标准正交基 $\{w_i\}$, 使得
其中至多一个(不妨设为 $w_1$)在 $x$ 处不为零. 利用迹在正交基变换下的
不变性, 有
\begin{equation}
K(x)=\sum_{i=1}^{\mathcal{N}}|w_i|^2(x)=w_1^2(x).
\label{eq:3.92}
\end{equation}
现在, 由于每个 $w_i$ 在 $B_r$ 上的 $L^2$ 范数为 $1$, 平均值不等式给出:
对 $0<s\leq\frac12$ 与任意 $x\in B_{(1-s)r}$ 有
\begin{equation}
|w_i|^2(x)\leq\frac{1}{\Vol(B_{sr}\subset\R^k)}
 \int_{B_{sr}(x)}w_i^2
 \leq\frac{\mathrm{V}_\Sigma\,s^{-k}}{\Vol(B_r\cap\Sigma)}.
\label{eq:3.93}
\end{equation}
结合 \eqref{eq:3.92} 与 \eqref{eq:3.93}, 得到: 对每个 $j=1,\dots,N$ 与任意
$x\in B_{(1-j/(2d))r}$ 有
\begin{equation}
\Vol(B_r\cap\Sigma)K(x)
 =\Vol(B_r\cap\Sigma)|w_1|^2(x)
 \leq\mathrm{V}_\Sigma\Bigl(\frac{j}{2d}\Bigr)^{-k}.
\label{eq:3.94}
\end{equation}

我们把 $K$ 的积分拆开:
\begin{equation}
\int_{B_{(1-\frac{1}{2d})r}}K
 =\int_{B_{(1-\frac{N}{2d})r}}K
 +\sum_{j=1}^{N-1}\int_{B_{(1-\frac{j}{2d})r}\setminus
 B_{(1-\frac{j+1}{2d})r}}K.
\label{eq:3.95}
\end{equation}
现在估计 \eqref{eq:3.95} 右端的两个项. 由于 $d\leq N\leq2d$, 有
$B_{(1-N/(2d))r}\subset B_{\frac r2}$, 故由 \eqref{eq:3.94}
\begin{equation}
\int_{B_{(1-\frac{N}{2d})r}}K
 \leq\mathrm{V}_\Sigma\Bigl(\frac12\Bigr)^{-k}
 \frac{\Vol(B_{\frac r2}\cap\Sigma)}{\Vol(B_r\cap\Sigma)}
 \leq\mathrm{V}_\Sigma\,2^k.
\label{eq:3.96}
\end{equation}
用 $K$ 在每个环域上的最大值来估计其积分, \eqref{eq:3.94} 给出
\begin{equation}
\begin{aligned}
&\sum_{j=1}^{N-1}\int_{B_{(1-\frac{j}{2d})r}\setminus
 B_{(1-\frac{j+1}{2d})r}}K\\
&\quad\leq V_\Sigma\sum_{j=1}^{N-1}
 \frac{\Vol\bigl(B_{(1-\frac{j}{2d})r}\setminus
 B_{(1-\frac{j+1}{2d})r}\cap\Sigma\bigr)}{\Vol(B_r\cap\Sigma)}
 \Bigl(\frac{j}{2d}\Bigr)^{-k}\\
&\quad\leq V_\Sigma\sum_{j=1}^{N-1}2(k+1)
 \Bigl(\frac{1}{2d}\Bigr)\Bigl(\frac{j}{2d}\Bigr)^{-k},
\end{aligned}
\label{eq:3.97}
\end{equation}
其中第二个不等式由 \eqref{eq:3.87}(体积正则性)得到. 利用初等不等式
\begin{equation}
\sum_{j=1}^{N-1}j^{-k}=1+\sum_{j=2}^{N-1}j^{-k}
 \leq1+\int_1^\infty s^{-k}\,ds=\frac{k}{k-1},
\label{eq:3.98}
\end{equation}
由 \eqref{eq:3.97} 可得
\begin{equation}
\sum_{j=1}^{N-1}\int_{B_{(1-\frac{j}{2d})r}\setminus
 B_{(1-\frac{j+1}{2d})r}}K
 \leq V_\Sigma(k+1)\frac{k}{k-1}2^k d^{k-1}.
\label{eq:3.99}
\end{equation}
把 \eqref{eq:3.96} 与 \eqref{eq:3.99} 代入 \eqref{eq:3.95} 得
\begin{equation}
\begin{aligned}
\int_{B_{(1-\frac{1}{2d})r}}K
&\leq\mathrm{V}_\Sigma\,2^k
 \Bigl(1+(k+1)\frac{k}{k-1}d^{k-1}\Bigr)\\
&\leq V_\Sigma\,2^{k+1}(k+1)\frac{k}{k-1}d^{k-1},
\end{aligned}
\label{eq:3.100}
\end{equation}
因为 $d\geq1$ 且 $k>1$.

结合 \eqref{eq:3.88} 与 \eqref{eq:3.100}, 得到
\begin{equation}
\mathcal{N}\leq a^{-1}\int_{B_{(1-\frac{1}{2d})r}}K
 \leq\mathrm{V}_\Sigma(k+1)\frac{k}{k-1}a^{-1}2^{k+1}d^{k-1}.
\label{eq:3.101}
\end{equation}
\end{proof}

\thmref{thm:3.18} 现在可直接推出.

\begin{proof}[\thmref{thm:3.18} 的证明]
取 $\Omega=(1-1/(2d))^{-1}$, 并选取正整数 $m_0$, 使得
$\Omega^{m_0}\geq R_0$. 设 $u_1,\dots,u_{2s}\in\mathcal{H}_d(\Sigma)$
线性无关. 由\propref{prop:3.22}, 存在 $N\geq m_0$、整数 $\ell$ 满足
\begin{equation}
\frac{\mathrm{e}^{-4}}{2}s\leq\ell,
\label{eq:3.102}
\end{equation}
以及 $u_i$ 张成的线性空间中的 $J_r$-标准正交函数
$f_1,\dots,f_\ell$, 使得对 $i=1,\dots,\ell$ 有
\begin{equation}
\frac{\mathrm{e}^{-4}}{2}\leq I_{f_i}\bigl((1-(2d)^{-1})r\bigr),
\label{eq:3.103}
\end{equation}
其中我们记 $r=\Omega^{N+1}$. \propref{prop:3.23}(取 $a=\mathrm{e}^{-4}/2$)蕴含
\begin{equation}
\ell\leq\mathrm{V}_\Sigma(k+1)\frac{k}{k-1}
 \mathrm{e}^4\,2^{k+2}d^{k-1}.
\label{eq:3.104}
\end{equation}
结合 \eqref{eq:3.102} 与 \eqref{eq:3.104} 得
\begin{equation}
\dim\mathcal{H}_d(\Sigma)\leq\mathrm{V}_\Sigma(k+1)\frac{k}{k-1}
 \mathrm{e}^8\,2^{k+4}d^{k-1}.
\label{eq:3.105}
\end{equation}
\end{proof}

我们将以一段非常简短的讨论结束本章, 把本章的问题与欧氏空间中极小锥的
结果联系起来.

下述初等引理的证明留给读者(进一步的讨论可参看例如 \cite{CM11} 与
\cite{CM15}):

\begin{lemma}[参看第 2 章的 \eqref{eq:2.179}]
\label{lem:3.24}
若 $g$ 是 $N^{k-1}$ 上特征值为 $\lambda$ 的特征函数, 且
$p^2+(k-2)p=\lambda$, 则 $r^pg$ 是 $C(N)$ 上的调和函数. 事实上
$r^pg\in\mathcal{H}_p(C(N))$.
\end{lemma}

由这个引理可见, $N$ 的谱性质等价于锥 $C(N)$ 上多项式增长的调和函数的
性质. 从这个观点看, \thmref{thm:3.18} 大致是 Weyl 渐近公式的类比. 关于这一观点
在相关框架下的进一步发展, 见 \cite{CM15}.

球面极小子流形的谱性质本身也已被研究(例如见 Cheng、Li 与 Yau
\cite{CgLiYa}, Choi 与 Wang \cite{CiWa}(\thmref{thm:7.11}), 以及 Li 与 Tian
\cite{LiTi}).

% ================================================================
\section{气泡收敛蕴含变分流形收敛}
\label{sec:3.6}

本节中我们回顾从 $\bS^2$ 出发的一列映射的\newterm{bubble-convergence}{气泡收敛}{bubble convergence}概念, 然后证明气泡收敛
蕴含变分流形收敛. 这在 \cite{CM27} 中已证明. 这一事实以后在我们用
min-max 论证证明极小二维球面的存在性时会用到.

\mtpara{气泡收敛}
我们需要 $W^{1,2}$ 映射序列 $v^j$ 收敛到一族 $W^{1,2}$ 映射
$\{u_0,\dots,u_m\}$ 的某个概念, 其精神类似于 Gromov 关于伪全纯曲线的
紧性定理 \cite{G} 中的收敛. 我们要用的概念是 Parker 与 Wolfson 在
\cite{PkW} 中为 $J$-全纯曲线、Parker 在 \cite{Pk} 中为调和映射建立的
气泡树收敛的一个轻微减弱. 在我们的应用中, $v^j$ 近似调和, 而极限映射
$u_i$ 调和. 为把这一点说清楚, 我们需要下一个定义.

$S^+$ 与 $S^-$ 表示 $\bS^2$ 中的北半球与南半球, $p^+=(0,0,1)$ 与
$p^-=(0,0,-1)$ 表示北极与南极.

\begin{definition}
\label{def:3.25}
给定球 $B_r(x)\subset\bS^2$, 把 $B_r(x)$ 映到 $S^-$ 的\textbf{共形
伸缩}是指先作平移 $x\to p^-$, 再关于 $p^-$ 作把 $B_r(p^-)$ 映到 $S^-$
的 $\bS^2$ 的伸缩所得的复合.
\end{definition}

共形伸缩的标准例子来自: 先作球极投影
\begin{equation}
\Pi:\bS^2\setminus\{(0,0,1)\}\to\R^2,
\label{eq:3.106}
\end{equation}
再把 $\R^2$ 按正的 $\lambda\neq1$ 伸缩, 最后作 $\Pi^{-1}$.

在下面关于收敛的定义中, 映射 $u_0$ 将是 $v^j$ 的通常 $W^{1,2}$-弱极限
(见 (B1)), 而其余的 $u_i$ 将作为 $v^j$ 与一列发散的 $\bS^2$ 共形伸缩
之复合的弱极限出现(见 (B2)). 条件 (B3) 保证这些极限都是以真正不同的
方式出现的, 而条件 (B4) 意味着这些 $u_i$ 合起来占尽了全部能量.

\begin{definition}[气泡收敛]
\label{def:3.26}
若下列条件成立, 则称 $W^{1,2}$ 映射序列 $v^j:\bS^2\to M$ 收敛到一族
$W^{1,2}$ 映射 $u_0,\dots,u_m:\bS^2\to M$:
\begin{itemize}[leftmargin=3.6em,labelwidth=3.4em,align=left,itemsep=3pt]
\item[\textbf{(B1)}] $v^j$ 在 $W^{1,2}$ 中弱收敛于 $u_0$, 且存在有限集
 $S_0=\{x_0^1,\dots,x_0^{k_0}\}\subset\bS^2$, 使得对任意紧集
 $K\subset\bS^2\setminus S_0$, $v^j$ 在 $W^{1,2}(K)$ 中强收敛于 $u_0$.
\item[\textbf{(B2)}] 对每个 $i>0$, 我们得到一个点
 $x_{\ell_i}\in S_0$ 与一列球 $B_{r_{i,j}}(y_{i,j})$, 其中
 $y_{i,j}\to x_{\ell_i}$, $r_{i,j}\to0$. 此外, 若
 $D_{i,j}:\bS^2\to\bS^2$ 是把南半球映到
 $B_{r_{i,j}}(y_{i,j})$ 的共形伸缩, 则映射 $v^j\circ D_{i,j}$ 如
 (B1) 那样收敛到 $u_i$. 即: $v^j\circ D_{i,j}$ 在 $W^{1,2}(\bS^2)$ 中
 弱收敛于 $u_i$, 且存在有限集 $S_i$, 使得对任意紧集
 $K\subset\bS^2\setminus S_i$, $v^j\circ D_{i,j}$ 在
 $W^{1,2}(K)$ 中强收敛.
\item[\textbf{(B3)}] 若 $i_1\neq i_2$, 则
 $\displaystyle\frac{r_{i_1,j}}{r_{i_2,j}}+\frac{r_{i_2,j}}{r_{i_1,j}}
 +\frac{|y_{i_1,j}-y_{i_2,j}|^2}{r_{i_1,j}r_{i_2,j}}\to\infty$.
\item[\textbf{(B4)}] 能量等式成立:
 $\displaystyle\sum_{i=0}^mE(u_i)=\lim_{j\to\infty}E(v^j)$.
\end{itemize}
\end{definition}

\mtpara{气泡收敛的两个简单例子}
气泡收敛最简单的非平凡例子是: 每个映射 $v^j=u\circ\Psi_j$ 都是固定的
调和映射 $u:\bS^2\to M$ 与一列发散的伸缩 $\Psi_j:\bS^2\to\bS^2$ 的
复合. 此时 $v^j$ 在 $\bS^2\setminus\{p_-\}$ 的每个紧集上收敛于常值映射
$u_0=u(p_+)$, 而全部能量集中在单点 $p_-=S_0$ 处. 把 $v^j$ 与发散的
共形伸缩序列 $\Psi_j^{-1}$ 复合, 就得到极限 $u_1=u$.

第二个例子如下. 设
\begin{equation}
\Pi:\bS^2\setminus\{(0,0,1)\}\to\R^2
\label{eq:3.107}
\end{equation}
是球极投影, $z=x+\ii y$ 是 $\R^2=\C$ 上的复坐标. 若令
\begin{equation}
f_j(z)=1/(jz)+z=\frac{z^2+1/j}{z},
\label{eq:3.108}
\end{equation}
则映射
\begin{equation}
v^j=\Pi^{-1}\circ f_j\circ\Pi:\bS^2\to\bS^2
\label{eq:3.109}
\end{equation}
是共形的, 从而是调和的. 由于每个 $v^j$ 是二次有理映射, 有
$E(v^j)=\Area(v^j)=8\pi$. 此外, $v^j$ 在 $0$ 之外收敛于恒等映射, 后者
的能量为 $4\pi$. 另外 $4\pi$ 的能量在 $0$ 处消失, 但可以通过与一列发散
的共形伸缩复合而由一个映射 $u_1$ 来解释; $u_1$ 的次数也必须为 $1$. 此时, 这些
共形伸缩把 $f_j$ 变为
\begin{equation}
\tilde{f}_j(z)=f_j(z/j)=1/z+z/j,
\label{eq:3.110}
\end{equation}
它收敛于关于半径为一的圆的反演(共形).

\mtpara{气泡收敛蕴含变分流形收敛}

\begin{proposition}[Colding--Minicozzi, \cite{CM27}]
\label{prop:3.27}
若 $W^{1,2}(\bS^2,M)$ 映射序列 $v^j$ 气泡收敛到有限个光滑映射
$u_0,\dots,u_m:\bS^2\to M$, 则它也变分流形收敛.
\end{proposition}

在进入证明之前, 回忆: 若对一切 $\delta>0$, 集合
\begin{equation}
\{x\mid|f_j-f|(x)>\delta\}
\label{eq:3.111}
\end{equation}
的测度当 $j\to\infty$ 时趋于零, 则称函数序列 $f_j$ 依测度收敛于 $f$;
见 \cite{R} 第 95 页. 显然, $L^1$ 收敛蕴含依测度收敛. 此外, 若 $f_j$
依测度收敛于 $f$, 且 $h$ 一致连续, 则
\[
h\circ f_j\to h\circ f
\]
依测度收敛. 最后, 我们将用到下述一般形式的控制收敛定理, 它综合了
\cite{R} 第 92 页定理 17 与第 96 页命题 20:

\smallskip
\noindent\textbf{(DCT)} 若 $f_j\to f$ 依测度收敛, $g_j\to g$ 在 $L^1$ 中
收敛, 且 $|f_j|\leq g_j$, 则 $\int f_j\to\int f$.
\smallskip

我们还将用到: 映射 $\nabla u\mapsto J_u$ 作为从 $L^2$ 到 $L^1$ 的映射是
连续的, 从而 $\Area(u)$ 关于 $E(u)$ 连续. 确切地说, 若
$u,v\in W^{1,2}(\bS^2,M)$, 则
\begin{equation}
|J_u-J_v|\leq\sqrt2\,|\nabla u-\nabla v|^{1/2}
 \max\bigl\{|\nabla u|^{3/2},|\nabla v|^{3/2}\bigr\}.
\label{eq:3.112}
\end{equation}
这由下述线性代数事实推出\footnote{注意 $|S^TT|\leq|S||T|$,
$|\operatorname{Tr}(S^TT)|\leq|S||T|$, 且若 $X_t$ 是一族 $2\times2$ 矩阵,
则 $\partial_t\det X_t=\operatorname{Tr}(X_t^c\partial_tX_t)$, 其中
$X_t^c$ 是把对角元互换、非对角元乘以 $-1$ 所得之余子矩阵. 把它应用于
$X_t=(S+t(T-S))^T(S+t(T-S))$, 并利用中值定理, 即得 \eqref{eq:3.113}.}:
若 $S$ 与 $T$ 是 $N\times2$ 矩阵, 则
\begin{equation}
\bigl|\det(S^TS)-\det(T^TT)\bigr|
 \leq2|T-S|\max\{|S|^3,|T|^3\},
\label{eq:3.113}
\end{equation}
其中 $|S|^2$ 是 $S$ 的各元素平方之和, $S^T$ 是转置.

\begin{proof}[\propref{prop:3.27} 的证明]
对每个 $v^j$, 记 $V^j$ 为相应的到 $G_2M$ 的映射. 同样, 对每个 $u_i$,
记 $U_i$ 为相应的到 $G_2M$ 的映射.

由 (B1)--(B4), 我们可以选取 $m+1$ 列区域
$\Omega_0^j,\dots,\Omega_m^j\subset\bS^2$, 使得对每个 $j$ 它们两两不交,
并且对每个 $i=0,\dots,m$, 把 $D_{i,j}^{-1}$ 作用于 $\Omega_i^j$ 得到
一列收敛到 $\bS^2\setminus S_i$ 的区域, 且占尽全部能量, 即
\begin{equation}
\lim_{j\to\infty}\int_{\bS^2\setminus(\bigcup_i\Omega_i^j)}
 |\nabla v^j|^2=0.
\label{eq:3.114}
\end{equation}
由 \eqref{eq:3.114}, 命题归结为: 对每个 $i$ 与任意 $h\in C^0(G_2M)$, 证明
\begin{equation}
\begin{aligned}
\int_{\bS^2}h\circ U_i\,J_{u_i}
&=\lim_{j\to\infty}\int_{\Omega_i^j}h\circ V^j\,J_{v^j}\\
&=\lim_{j\to\infty}\int_{D_{i,j}^{-1}(\Omega_i^j)}
 h\circ V^j\circ D_{i,j}\,J_{(v^j\circ D_{i,j})},
\end{aligned}
\label{eq:3.115}
\end{equation}
其中最后一个等号只是积分的变量替换公式.

为在 \eqref{eq:3.115} 的证明中简化记号, 对每个 $i$ 与 $j$, 记 $v_i^j$ 为
$v^j\circ D_{i,j}$ 在 $D_{i,j}^{-1}(\Omega_i^j)$ 上的限制, 记 $V_i^j$
为相应的到 $G_2M$ 的映射.

注意由 \eqref{eq:3.112}, $J_{v_i^j}\to J_{u_i}$ 在 $L^1(\bS^2)$ 中成立. 给定
$\epsilon>0$ 与 $i$, 定义
\[
\Omega_\epsilon^i=\bigl\{J_{u_i}\geq\epsilon\bigr\}.
\]
由于 $h$ 有界且 $J_{v_i^j}\to J_{u_i}$ 于 $L^1(\bS^2)$, 只需证明
\begin{equation}
\lim_{j\to\infty}\int_{\Omega_\epsilon^i}h\circ V_i^j\,J_{v_i^j}
 =\int_{\Omega_\epsilon^i}h\circ U_i\,J_{u_i},
\label{eq:3.116}
\end{equation}
\eqref{eq:3.115} 便随之成立. 然而, 对任意 $\delta>0$, $W^{1,2}$ 收敛蕴含集合
\begin{equation}
\bigl\{x\in\Omega_\epsilon^i\mid J_{v_i^j}\geq\frac\epsilon2
 \text{ 且 } |V_i^j-U_i|\geq\delta\bigr\}
\label{eq:3.117}
\end{equation}
的测度当 $j\to\infty$ 时趋于零. 由于 Jacobi 因子的 $L^1$ 收敛蕴含集合
$\{x\in\Omega_\epsilon^i\mid J_{v_i^j}<\frac\epsilon2\}$ 的测度趋于零,
可知映射 $V_i^j$ 在 $\Omega_\epsilon^i$ 上依测度收敛于 $U_i$. 因此
$h\circ V_i^j$ 在 $\Omega_\epsilon^i$ 上依测度收敛于 $h\circ U_i$. 于是
一般形式的控制收敛定理 (DCT) 给出 \eqref{eq:3.116}, 从而也给出 \eqref{eq:3.115}.
\end{proof}
